%%%%%%%%%%%%%%%%%%%%%%%%%%%%%%%%%%%%%%%%%%%%%%%%
% BOZZA PER TESI DI LAUREA IN LATEX
%%%%%%%%%%%%%%%%%%%%%%%%%%%%%%%%%%%%%%%%%%%%%%%%

\documentclass[12pt,titlepage]{book}
\usepackage[italian]{babel}
%\usepackage{graphics}
\usepackage{url,amsfonts,epsfig}
\usepackage[applemac]{inputenc} %comando per le lettere accentate se usate mac  
%\usepackage[latin1]{inputenc} % comando per le lettere accentate se usate pc  


\title{\textsc{Questo � il titolo della mia tesi}}
\author{Pinco Pallino}

\begin{document}
\pagenumbering{roman}

%%%% Opzione per interlinea 2
%%%\baselineskip 18pt

\maketitle

\tableofcontents
%%\listoffigures
%%\listoftables

\chapter{Introduzione} \label{introduzione}
\pagenumbering{arabic}

Questo \`e un file minimale d'esempio per la
stesura di tesi in \LaTeX.
Ricordatevi di dare il comando bibtex per generare la
bibliografia (gli items bibliografici vanno messi
nel file *.bib).\footnote{Ricordatevi di usare
 meno footnotes (note) possibili.}
Come esempio, caratteri matematici si possono inserire
facilmente nel testo: $y_{it} = \alpha+\vec{x}_{it}^{T}\beta_{it} +\epsilon_{it} $ ma
anche centrati:
$$ y_{it} = \alpha+\vec{x}_{it}^{T}\beta_{it} +\epsilon_{it}
$$ con $$ \epsilon_{it}\sim N(0,\sigma^2)
$$



Possiamo facilmente citare altri lavori
usando le loro labels~\cite{dominici03,Deb97}.\\
Possiamo creare tabelle:
\begin{center}
\begin{tabular}{|c|p{8cm}|}
\hline
Variabile & Descrizione \\
\hline \hline
\textit{Income} - $Y_i$ & reddito individuale mensile in dollari per soggetto i-esimo\\
\hline
\textit{Education} - $X_i$ & livello d'istruzione del soggetto i-esimo\\
\hline\hline
\end{tabular}
\end{center}
Le tabelle potete metterle in appositi ambienti
che il compilatore decider� dove inserire,
\begin{table}[ht]
\begin{center}
\begin{tabular}{|l|c|c|}
\hline
{\bf Lag lenght} ($p$) &
{\bf Personal income} ($Y$) & {\bf Change in personal income} ($\Delta Y$)\\\hline
{\bf 1}  & 0.999 & -0.010 \\
{\bf 2} & 0.999 & 0.012\\
{\bf 3} & 0.999 & 0.134\\
{\bf 4} & 0.998 & 0.008\\
{\bf 5} & 0.998 & -0.15 \\
\hline
\end{tabular}
\end{center}
\caption{Funzioni di autocorrelazione
\label{tab:autocorrelation}}
\end{table}
come ad esempio in Tabella~\ref{tab:autocorrelation}.

Si possono inoltre importare disegni in vari formati,
tra cui l'eps (encapsulated postscript) garantisce un ottimo
risultato. Ad esempio, con il comando
\begin{verbatim}
\psfig{figure=palline.eps,width=0.8\textwidth}
\end{verbatim}
si include la figura dal nome {\tt palline.eps}
e si fissa la sua dimensione orizzontale all'80\% della larghezza del foglio. La dimensione verticale
sar� proporzionale a seconda del disegno. Si pu� forzare anche quella
aggiungendo dentro la parentesi grafa l'opportuno comando, ad esempio: {\tt, height=3cm}.
%In Figura~\ref{includieps} si pu� vedere l'effetto dell'istruzione suddetta:

%\begin{figure}[ht]
%\begin{center}
%\psfig{figure=palline.eps,width=0.8\textwidth}
%\end{center}
%\caption{Inclusione di figura esterna\label{includieps}}
%\end{figure}

\medskip


La conclusione dell'Introduzione
finir� pi� o meno in questo modo:


\medskip

La struttura della tesi
\`{e} la seguente.
Nel Capitolo~\ref{cap:econtheory},
bla bla ...............

Infine, nel Capitolo~\ref{cap:conclusioni},
.......

In Appendice sono riportate le principali procedure
del codice sviluppato.



\chapter{La teoria economica su ....} \label{cap:econtheory}
\begin{flushright}
\textit{``In economia, gran parte della saggezza\\ consiste nel sapere ci� che non sai''.\\
John Kenneth Galbraith}
\end{flushright}
\medskip
\section{Il modello di...}\label{sec:mod1}
Bla, bla .... come vedremo poi nella sezione~\ref{sec:crit} ... bla bla.


\section{Critiche al modello di ....}\label{sec:crit}

Riprendiamo quanto visto nella sezione~\ref{sec:mod1}, bla bla...



\chapter{Conclusioni}
\label{cap:conclusioni}
Nella tesi abbiamo dunque mostrato che .......



Ultima nota. Se la tesi vi sembrasse corta, potreste
aumentare l'interlinea col comando {\tt baselineskip}.
Guardate qui l'effetto e poi guardate il sorgente per vedere
come si usa.

\baselineskip 25pt

Sulla base di queste premesse, Wagstaff, definisce una funzione di utilit� chiamata $U$, definita in modo tale che sia concava.....essa pu� essere denotata da $H(M)$, tale per cui $\partial H(M)/ \partial M>0$ con $\partial^{2} H(M)/\partial M^{2}<0$. 

\baselineskip 5pt

Di tutte le variabili che possono essere incluse nel nostro fattore $M$, quelle che hanno maggiore influenza nell'andamento di $M$ e quindi di $H$, sono la conoscenza applicata ...quindi la funzione $H(M)$.

\baselineskip 16pt


%%% EVENTUALE
\appendix
\chapter{Codice}

\begin{verbatim}
num(0).
num(s(X)) :- num(X).
\end{verbatim}

%%% OBBLIGATORIA:

\bibliographystyle{plain}
\bibliography{biblio} %%% nome file(s)

\end{document}